\part{Reference}

% =====================================================================
\chapter{Every tag in CompLaB.xml}
\label{ch:reference}

This chapter lists every element the code reads, with its type, default and
units. Where a tag is mandatory only under some condition, that condition is
stated: those are the ones that stop a run.

\note{The file that documents itself.}{\file{CompLaB\_reference\_template.xml}
ships with the code and carries all of this as inline comments, next to a
plausible value. When you are actually editing a file, that is faster than this
chapter.}

\section{\tg{path}}

\begin{longtable}{@{}p{3.2cm}p{1.6cm}p{2.0cm}p{7.6cm}@{}}
\toprule
\textbf{Tag} & \textbf{Type} & \textbf{Default} & \textbf{Meaning}\\
\midrule\endhead
\tg{src\_path} & string & \code{src} & where \file{complab3d\_cobrapy.py} is found at run time. Only used by the COBRApy path, but it is imported by name, so the file must be there.\\
\tg{input\_path} & string & \code{input} & geometry and metabolic model files\\
\tg{output\_path} & string & \code{output} & VTK and checkpoint files. Must exist; the code does not create it.\\
\bottomrule
\end{longtable}

\section{\tg{simulation\_mode}}

\begin{longtable}{@{}p{4.6cm}p{1.4cm}p{1.6cm}p{6.8cm}@{}}
\toprule
\textbf{Tag} & \textbf{Type} & \textbf{Default} & \textbf{Meaning}\\
\midrule\endhead
\tg{biotic\_mode} & bool & \code{true} & \code{false} skips the whole \tg{microbiology} block and runs pure chemistry\\
\tg{enable\_kinetics} & bool & \code{true} & run \file{defineKinetics.hh}\\
\tg{enable\_abiotic\_kinetics} & bool & \code{false} & run \file{defineAbioticKinetics.hh}\\
\tg{enable\_fba\_glpk} & bool & \code{false} & global switch for GLPK. Needs \code{-DENABLE\_GLPK=ON}.\\
\tg{enable\_fba\_cobrapy} & bool & \code{false} & global switch for COBRApy. Needs \code{-DENABLE\_COBRAPY=ON}.\\
\tg{enable\_surrogate} & bool & \code{false} & global switch for the surrogate network\\
\tg{fba\_concentration\_basis} & string & \code{free} & \code{free} or \code{total}. With speciation on, whether the uptake bound is built from the free ion or the total.\\
\tg{glpk\_lpsolver} & int & 1 & 1 simplex, 2 interior point. GLPK only.\\
\tg{glpk\_save\_problem} & bool & \code{false} & dump the LP of the first solve, for debugging\\
\tg{enable\_validation\_diagnostics} & bool & \code{false} & extra consistency printing at start-up. Worth switching on while setting a case up.\\
\bottomrule
\end{longtable}

\warn{The two gates must agree.}{A microbe asking for \code{glpk} while
\tg{enable\_fba\_glpk} is \code{false} terminates the run with a message. So
does a switch that is on in an executable built without the dependency. That is
deliberate: you cannot silently get different biology from the one you asked
for.}

\section{\tg{LB\_numerics}}

\subsection{\tg{domain}}

\begin{longtable}{@{}p{4.2cm}p{1.6cm}p{1.6cm}p{7.0cm}@{}}
\toprule
\textbf{Tag} & \textbf{Type} & \textbf{Default} & \textbf{Meaning}\\
\midrule\endhead
\tg{nx} \tg{ny} \tg{nz} & int & --- & domain size in voxels. Mandatory.\\
\tg{dx} & float & --- & voxel size, in \tg{unit}\\
\tg{unit} & string & \code{um} & \code{m}, \code{mm} or \code{um}\\
\tg{characteristic\_length} & float & \code{ny} & length scale for the P\'eclet and Reynolds numbers, in voxels\\
\tg{filename} & string & --- & geometry file, inside \tg{input\_path}\\
\bottomrule
\end{longtable}

\subsection{\tg{material\_numbers}}

\begin{longtable}{@{}p{3.4cm}p{2.4cm}p{1.4cm}p{7.6cm}@{}}
\toprule
\textbf{Tag} & \textbf{Type} & \textbf{Default} & \textbf{Meaning}\\
\midrule\endhead
\tg{pore} & int list & 2 & open fluid. A list, so several codes can be pore.\\
\tg{solid} & int & 0 & no dynamics at all\\
\tg{bounce\_back} & int & 1 & no-slip wall\\
\tg{microbe}$n$ & int list & --- & voxels where microbe $n$ is seeded. Its presence also marks the organism as a biofilm former.\\
\bottomrule
\end{longtable}

\subsection{Flow and time stepping}

\begin{longtable}{@{}p{4.0cm}p{1.4cm}p{1.8cm}p{7.2cm}@{}}
\toprule
\textbf{Tag} & \textbf{Type} & \textbf{Default} & \textbf{Meaning}\\
\midrule\endhead
\tg{delta\_P} & float & 0 & pressure drop driving the flow, lattice units\\
\tg{Peclet} & float & 0 & 0 means no flow: pure diffusion, and the Navier--Stokes solver never runs\\
\tg{tau} & float & 0.8 & lattice-Boltzmann relaxation time. Values below about 0.55 are unstable.\\
\tg{track\_performance} & bool & \code{false} & timing output. Suppresses all VTI and checkpoint writing.\\
\tg{ns\_max\_iT1} & int & 100000 & flow solver step cap, first solve\\
\tg{ns\_max\_iT2} & int & 100000 & the same for later re-solves, after the geometry changes\\
\tg{ns\_converge\_iT1} & float & 1e-8 & ``flow is steady'' tolerance, first solve\\
\tg{ns\_converge\_iT2} & float & 1e-6 & the same, later\\
\tg{ns\_update\_interval} & int & 1 & re-solve the flow every $N$ reaction steps\\
\tg{ade\_update\_interval} & int & 1 & apply reactions every $N$ transport steps\\
\tg{ade\_max\_iT} & int & 1e7 & reaction and transport steps to run\\
\tg{ade\_converge\_iT} & float & 1e-8 & stop early if the chemistry stops changing\\
\bottomrule
\end{longtable}

\section{\tg{chemistry}}

\begin{longtable}{@{}p{4.6cm}p{1.8cm}p{1.4cm}p{6.6cm}@{}}
\toprule
\textbf{Tag} & \textbf{Type} & \textbf{Default} & \textbf{Meaning}\\
\midrule\endhead
\tg{number\_of\_substrates} & int & --- & must equal the number of \tg{substrateN} blocks\\
\tg{fix\_concentration} & bool list & all false & one per substrate. Marks an inexhaustible reservoir: read but never drawn down. Accepts \code{yes}/\code{no} and \code{1}/\code{0}.\\
\tg{fix\_lower\_bounds} & bool list & all false & one per substrate. Clamps the FBA uptake bound to what is locally present.\\
\bottomrule
\end{longtable}

\subsection{Each \tg{substrateN}}

\begin{longtable}{@{}p{5.4cm}p{1.4cm}p{1.4cm}p{6.2cm}@{}}
\toprule
\textbf{Tag} & \textbf{Type} & \textbf{Default} & \textbf{Meaning}\\
\midrule\endhead
\tg{name\_of\_substrates} & string & --- & used in output and to match \tg{components}\\
\tg{initial\_concentration} & float & 0 & mol/L everywhere at $t=0$\\
\tg{substrate\_diffusion\_} \tg{coefficients}/\tg{in\_pore} & float & 1e-9 & m$^2$/s in open water\\
\quad \dots/\tg{in\_biofilm} & float & 2e-10 & m$^2$/s inside biofilm\\
\tg{immobile} & bool & \code{false} & \code{true} means never advects or diffuses; receives only its reaction source. How minerals are represented.\\
\tg{left\_boundary\_type} & string & --- & \code{Dirichlet} or \code{Neumann}\\
\tg{left\_boundary\_condition} & float & 0 & the value: concentration for Dirichlet, flux for Neumann\\
\tg{right\_boundary\_type} & string & --- & as above\\
\tg{right\_boundary\_condition} & float & 0 & as above\\
\bottomrule
\end{longtable}

\section{\tg{microbiology}}

\begin{longtable}{@{}p{4.6cm}p{1.4cm}p{1.6cm}p{6.8cm}@{}}
\toprule
\textbf{Tag} & \textbf{Type} & \textbf{Default} & \textbf{Meaning}\\
\midrule\endhead
\tg{number\_of\_microbes} & int & 0 & must equal the number of \tg{microbeN} blocks\\
\tg{maximum\_biomass\_density} & float & --- & gDW/L at which a voxel is full. Mandatory with CA.\\
\tg{thrd\_biofilm\_fraction} & float & --- & fraction above which a voxel counts as biofilm. Mandatory with CA.\\
\tg{CA\_method} & string & \code{fraction} & how the cellular automaton shares excess biomass\\
\bottomrule
\end{longtable}

\subsection{Each \tg{microbeN}: always}

\begin{longtable}{@{}p{5.4cm}p{1.4cm}p{1.4cm}p{6.2cm}@{}}
\toprule
\textbf{Tag} & \textbf{Type} & \textbf{Default} & \textbf{Meaning}\\
\midrule\endhead
\tg{name\_of\_microbes} & string & --- & used in output\\
\tg{solver\_type} & string & --- & \code{CA}, \code{FD} or \code{LBM}. Mandatory.\\
\tg{reaction\_type} & string & \code{kinetics} & see the table below\\
\tg{initial\_densities} & float list & --- & gDW/L. One entry per material number the organism is seeded in.\\
\tg{decay\_coefficient} & float & 0 & 1/s, first order\\
\tg{half\_saturation\_constants} & float list & --- & mol/L, one per substrate. Feeds the FBA and surrogate uptake bounds.\\
\tg{maximum\_uptake\_flux} & float list & all 0 & the \textbf{kinetics} path's Vmax, mol substrate per mol cells per second\\
\tg{left\_boundary\_type} \dots & --- & Neumann 0 & as for substrates\\
\bottomrule
\end{longtable}

\subsection{Each \tg{microbeN}: conditionally}

\begin{longtable}{@{}p{5.4cm}p{4.0cm}p{5.2cm}@{}}
\toprule
\textbf{Tag} & \textbf{Required when} & \textbf{Meaning}\\
\midrule\endhead
\tg{viscosity\_ratio\_in\_biofilm} & the organism has a \tg{microbeN} material entry & ratio used for the flow through biofilm. Must be present for \emph{exactly} the seeded organisms; the parser counts them.\\
\tg{biomass\_diffusion\_} \tg{coefficients} & \code{solver\_type} is FD & m$^2$/s, \tg{in\_pore} and \tg{in\_biofilm}\\
\tg{model\_filename} & \code{reaction\_type} uses FBA, unless \tg{model\_source} supplies it & the metabolic model, inside \tg{input\_path}. May be an SBML file (\file{iJO1366.xml}) or the \file{extractMM.py} matrix format; the root element decides which. A bare stem gets \file{.xml} appended, as it always did.\\
\tg{exchange\_reaction\_indices} & FBA, unless \tg{exchange\_reaction\_names} is given & one per substrate; \code{-1} (or \code{-99}) for one the organism does not exchange. Correct only for the exact model revision it was written against.\\
\tg{exchange\_reaction\_names} & alternative to the above; needs an SBML model & one reaction name per substrate, \code{none} for one the organism does not exchange. Resolved against the model at start-up, so a wrong name stops the run and suggests near matches. Wins if both are given.\\
\tg{fba\_maximum\_uptake\_flux} & FBA or surrogate & the \textbf{metabolic} path's Vmax, mmol/gDW/h. Separate from \tg{maximum\_uptake\_flux}.\\
\tg{substrate\_lower\_bounds} & FBA & hard floor on each exchange flux; eating is negative\\
\tg{substrate\_upper\_bounds} & FBA & hard ceiling\\
\tg{objective\_direction} & FBA & \code{maximize} or \code{minimize}, default maximize\\
\tg{biomass\_molar\_mass} & FBA or surrogate & gDW per mol of cells, default 24.6\\
\tg{equate\_bounds} & optional, FBA & \code{yes} forces uptake to exactly the bound instead of letting the solver choose anything up to it. Default \code{no}.\\
\tg{constraint\_indices} & optional, FBA & extra reactions to re-bound at load time\\
\tg{constraint\_lower\_bounds} & with the above & same length\\
\tg{constraint\_upper\_bounds} & with the above & same length\\
\bottomrule
\end{longtable}

\subsection{\tg{reaction\_type} spellings}

\begin{center}
\begin{tabular}{@{}p{6.4cm}p{8.4cm}@{}}
\toprule
\textbf{Accepted} & \textbf{Meaning}\\
\midrule
\code{none} \ \code{no} \ \code{0} & no metabolism\\
\code{kinetics} \ \code{kns} \ \code{1} & \file{defineKinetics.hh}\\
\code{glpk} \ \code{fba} \ \code{fba\_glpk} \ \code{2} & flux balance analysis through GLPK\\
\code{surrogate} \ \code{srg} \ \code{ann} \ \code{3} & the trained network\\
\code{cobrapy} \ \code{cpy} \ \code{fba\_cobrapy} \ \code{4} & flux balance analysis through COBRApy\\
\code{glpk\_and\_kinetics} \ \code{5} & both, rates added\\
\code{surrogate\_and\_kinetics} \ \code{6} & both\\
\code{cobrapy\_and\_kinetics} \ \code{7} & both\\
\bottomrule
\end{tabular}
\end{center}

\section{\tg{model\_source} and friends}

Top level, inside \tg{parameters}. Chapter \ref{ch:modelsource} explains the
search order.

\begin{longtable}{@{}p{3.6cm}p{1.4cm}p{2.2cm}p{7.2cm}@{}}
\toprule
\textbf{Tag} & \textbf{Type} & \textbf{Default} & \textbf{Meaning}\\
\midrule\endhead
\tg{model\_source} & string & --- & \code{bigg:<id>}, e.g. \code{bigg:iJO1366}. Absent means every FBA microbe supplies its own \tg{model\_filename}.\\
\tg{model\_bundle} & string & \code{models} & where the gzipped bundled models live, relative to the working directory\\
\tg{model\_cache} & string & \code{input} & where the unpacked \file{.xml} is put, and looked for first\\
\tg{allow\_download} & bool & \code{false} & permit fetching over the network when the model is neither cached nor bundled. Compute nodes usually cannot, which is why this is off.\\
\bottomrule
\end{longtable}

\note{Who gets the file.}{Every microbe whose \tg{reaction\_type} is metabolic
--- FBA \emph{or} surrogate --- and that did not name its own
\tg{model\_filename}. A surrogate microbe needs it too, because that is what
\tg{train\_if\_missing} sweeps.}

\section{Geometry generation, inside \tg{domain}}

All optional. Absent means \tg{filename} is read as before. Chapter
\ref{ch:geomgen} has the details.

\begin{longtable}{@{}p{3.4cm}p{1.4cm}p{1.8cm}p{7.6cm}@{}}
\toprule
\textbf{Tag} & \textbf{Type} & \textbf{Default} & \textbf{Meaning}\\
\midrule\endhead
\tg{generate} & string & --- & \code{channel}, \code{spheres}, \code{cylinders}, \code{random}, \code{fracture} or \code{layered}\\
\tg{porosity} & real & 0.5 & target open fraction, for \code{spheres}, \code{cylinders} and \code{random}\\
\tg{grain\_radius} & real & 3 & voxels, for \code{spheres} and \code{cylinders}\\
\tg{grain\_count} & int & --- & fix the number of grains instead of the porosity\\
\tg{aperture} & real & --- & voxels, for \code{fracture}\\
\tg{roughness} & real & 0 & fracture wall roughness, 0 to 1\\
\tg{layers} & int & --- & for \code{layered}\\
\tg{seed} & int & 1 & the generator is deterministic given this\\
\tg{walls} & string & --- & axes to wall off, e.g. \code{y} or \code{yz}\\
\tg{write\_geometry} & string & \code{generated\_geometry.dat} & where to record what was generated\\
\tg{import\_raw} & string & --- & a flat binary volume to threshold instead. Mutually exclusive with \tg{generate}.\\
\tg{raw\_dtype} & string & \code{uint8} & \code{uint8}, \code{uint16}, \code{int32}, \code{float32}, \code{float64}\\
\tg{threshold} & real & 128 & at or above this is one phase, below is the other\\
\tg{invert} & bool & \code{false} & swap which side of the threshold is pore\\
\bottomrule
\end{longtable}

\warn{The run stops on a sealed domain.}{Not percolating along $x$ is fatal, by
design. See chapter \ref{ch:geomgen}.}

\section{\tg{surrogate}}

Top level, inside \tg{parameters}. Chapter \ref{ch:surrogateblock} explains what
happens at start-up.

\begin{longtable}{@{}p{3.8cm}p{1.4cm}p{1.8cm}p{7.2cm}@{}}
\toprule
\textbf{Tag} & \textbf{Type} & \textbf{Default} & \textbf{Meaning}\\
\midrule\endhead
\tg{enabled} & bool & \code{false} & the whole block is ignored when this is off\\
\tg{weights\_file} & string & --- & a \file{.srg} file. Required when the block is enabled.\\
\tg{train\_if\_missing} & bool & \code{false} & fit one now if the file is not there, instead of stopping\\
\tg{microbe} & int & $-1$ & which microbe the network is for. $-1$ means all of them.\\
\bottomrule
\end{longtable}

Inside \tg{train}, all optional except the inputs:

\begin{longtable}{@{}p{3.8cm}p{1.4cm}p{1.8cm}p{7.2cm}@{}}
\toprule
\textbf{Tag} & \textbf{Type} & \textbf{Default} & \textbf{Meaning}\\
\midrule\endhead
\tg{input0} \ldots \tg{input7} & string & --- & \code{<substrate> <low> <high> [log]}. At least one when training. Three at most --- past that the samples needed grow faster than the fit improves.\\
\tg{samples} & int & 2000 & points in the sweep\\
\tg{layers} & ints & \code{10 10} & hidden layer sizes\\
\tg{restarts} & int & 5 & independent fits; the best is kept\\
\tg{epochs} & int & 4000 & Adam iterations per restart\\
\tg{verify\_points} & int & 512 & fresh points, re-solved by the LP, that the network never saw\\
\tg{seed} & int & 1 & the fit is reproducible given this, on every compiler\\
\tg{log\_output} & bool & \code{false} & fit $\log_{10}$ of the growth rate\\
\bottomrule
\end{longtable}

\warn{The training range is not advice.}{It is written into the \file{.srg}
file and enforced at run time: evaluations outside it are clamped to the edge
and counted, and a run that clamps more than 5\% of the time says so at the
end. Set the ranges from what the simulation will actually visit --- the upper
end is \tg{fba\_maximum\_uptake\_flux}, because no voxel can ever request
more.}

\section{\tg{diagnostics}}

Top level, inside \tg{parameters}. Chapter \ref{ch:diagnostics} has the
details.

\begin{longtable}{@{}p{3.8cm}p{1.4cm}p{2.0cm}p{7.0cm}@{}}
\toprule
\textbf{Tag} & \textbf{Type} & \textbf{Default} & \textbf{Meaning}\\
\midrule\endhead
\tg{enabled} & bool & \code{false} & the whole block is ignored when this is off\\
\tg{summary\_csv} & string & \code{summary.csv} & written inside \tg{output\_path} unless the path is absolute\\
\tg{interval} & int & 0 & steps between rows. 0 follows the VTI interval.\\
\tg{tolerance} & real & \code{1e-6} & relative drift above which a conserved sum is reported as failed\\
\tg{conserve} \ \tg{conserve1} \ldots \tg{conserve31} & string & --- & a sum of substrate names, \code{+} separated, that ought to be constant\\
\bottomrule
\end{longtable}

\section{\tg{equilibrium}}

\begin{longtable}{@{}p{3.6cm}p{2.6cm}p{8.4cm}@{}}
\toprule
\textbf{Tag} & \textbf{Type} & \textbf{Meaning}\\
\midrule\endhead
\tg{enabled} & bool & default \code{false}\\
\tg{components} & string list & master species. Names must match \tg{name\_of\_substrates}.\\
\tg{stoichiometry}/\tg{species}$n$ & float list & one row per substrate, one number per component. A row of zeros excludes the species.\\
\tg{logK}/\tg{species}$n$ & float & formation constant. Components take 0.\\
\bottomrule
\end{longtable}

\section{\tg{precipitation}}

\begin{longtable}{@{}p{3.8cm}p{1.6cm}p{1.4cm}p{7.6cm}@{}}
\toprule
\textbf{Tag} & \textbf{Type} & \textbf{Default} & \textbf{Meaning}\\
\midrule\endhead
\tg{enabled} & bool & \code{false} & \\
\tg{solid\_substrate} & int & --- & index of the mineral substrate. Must be \tg{immobile}.\\
\tg{max\_precipRho} & float & --- & mol/L of a full voxel; $1000/V_m$ with $V_m$ in cm$^3$/mol\\
\tg{surface\_only} & int & 1 & 1 grows only on wall-adjacent voxels (heterogeneous nucleation); 0 grows anywhere\\
\tg{perm\_ratio} & float & 0 & 0 makes a filled voxel an impermeable wall\\
\tg{update\_interval} & int & --- & steps between geometry re-checks. Each check that changes anything re-solves the flow.\\
\bottomrule
\end{longtable}

\section{\tg{dissolution}}

\begin{longtable}{@{}p{3.8cm}p{1.6cm}p{1.4cm}p{7.6cm}@{}}
\toprule
\textbf{Tag} & \textbf{Type} & \textbf{Default} & \textbf{Meaning}\\
\midrule\endhead
\tg{enabled} & bool & \code{false} & independent of \tg{precipitation}\\
\tg{reopen\_fraction} & float & 0.9 & a voxel reopens below this fraction of full. Hysteresis; see below.\\
\tg{surface\_only} & int & 1 & 1 dissolves only wetted voxels\\
\tg{update\_interval} & int & --- & steps between geometry re-checks\\
\tg{phase}$n$/\tg{name} & string & \code{phase}$n$ & label, for output\\
\tg{phase}$n$/\tg{material\_number} & int & --- & which mask code this phase occupies\\
\tg{phase}$n$/\tg{substrate} & int & --- & which \tg{immobile} substrate holds its inventory\\
\tg{phase}$n$/\tg{full\_density} & float & --- & mol/L when completely full\\
\tg{phase}$n$/\tg{initial\_fill} & float & 0 & mol/L at $t=0$. Non-zero for an original grain; 0 for a precipitate.\\
\tg{phase}$n$/\tg{is\_precipitate} & bool & \code{false} & marks the phase that \tg{precipitation} creates, so its voxels can later reopen\\
\bottomrule
\end{longtable}

\warn{Only declared phases dissolve.}{Anything without a \tg{phaseN} block can
never dissolve, and that is deliberate. If ``a solid voxel with no mineral left
becomes pore'' were the rule, an inert quartz grain --- which never held any
mineral --- is already below any threshold, and your grain pack would dissolve
on the first step.}

\section{\tg{IO}}

\begin{longtable}{@{}p{3.8cm}p{1.6cm}p{1.8cm}p{7.2cm}@{}}
\toprule
\textbf{Tag} & \textbf{Type} & \textbf{Default} & \textbf{Meaning}\\
\midrule\endhead
\tg{read\_NS\_file} & bool & \code{false} & restart the flow from a checkpoint\\
\tg{read\_ADE\_file} & bool & \code{false} & restart the chemistry from a checkpoint\\
\tg{ns\_filename} & string & \code{nsLattice} & \\
\tg{mask\_filename} & string & \code{maskLattice} & \\
\tg{subs\_filename} & string & \code{subsLattice} & index and iteration appended\\
\tg{bio\_filename} & string & \code{bioLattice} & index and iteration appended\\
\tg{save\_VTK\_interval} & int & 1000 & steps between snapshots\\
\tg{save\_CHK\_interval} & int & 1e6 & steps between restart points\\
\tg{debug\_updRxn} & int & 0 & 1 prints reaction-application debug lines\\
\bottomrule
\end{longtable}

% =====================================================================
\chapter{When something goes wrong}
\label{ch:errors}

\section{The run refuses to start}

CompLB3D validates its input before doing any work and stops with a named
message. That is by design: a run that starts and produces nonsense is worse
than one that refuses.

\begin{longtable}{@{}p{7.0cm}p{7.8cm}@{}}
\toprule
\textbf{Message contains} & \textbf{Cause and fix}\\
\midrule\endhead
\code{The length of \dots does not match} & A vector is the wrong length. The message names the tag. Almost always a substrate or microbe added without updating every per-substrate list.\\
\code{reaction\_type \dots is not defined} & A spelling the parser does not accept. See the table in Chapter~\ref{ch:reference}.\\
\code{asks for reaction\_type glpk but <enable\_fba\_glpk> is false} & The per-microbe choice and the global switch disagree. Set the switch, or change the microbe.\\
\code{is true but this executable was built without GLPK} & Rebuild with \code{-DENABLE\_GLPK=ON}.\\
\code{solver\_type \dots is not defined} & Not \code{FD}, \code{CA} or \code{LBM}.\\
\code{must be defined when solver\_type is Finite Difference} & FD needs \tg{biomass\_diffusion\_coefficients}.\\
\code{biofilm\_permeability\_ratio vector does not match} & \tg{viscosity\_ratio\_in\_biofilm} is present for a different set of organisms than the ones seeded under \tg{material\_numbers}. It must be exactly the seeded ones.\\
\code{<S> has \dots entries but nmet*nrxn =} & The metabolic model file is truncated or mis-generated. Re-run \file{extractMM.py}.\\
\code{<exchange\_reaction\_indices> entry \dots is outside} & An index past the end of that organism's model.\\
\code{must be defined when solver\_type is Cellular Automata} & CA needs \tg{thrd\_biofilm\_fraction}.\\
\bottomrule
\end{longtable}

\note{Catch these before you build.}{\file{examples/validateExamples.py} applies
most of these rules to an input file in about a second, without compiling
anything. Point it at a folder containing your \file{CompLaB.xml}.}

\section{The run starts but the answer is wrong}

\begin{longtable}{@{}p{5.0cm}p{9.8cm}@{}}
\toprule
\textbf{Symptom} & \textbf{Usual cause}\\
\midrule\endhead
\code{[NEG!]} warnings & A reaction removed more than was present in one step. Reduce \tg{ade\_update\_interval}, or your rate constant is too large --- check the per-second conversion.\\
nothing grows & Several possibilities, in order of likelihood: rate constants 86400 times too small (a per-day constant used directly); \tg{enable\_kinetics} false; the organism seeded on a material number that does not occur in the geometry; substrate never reaching it.\\
biomass explodes & No decay, no maximum density, or growth not limited by anything. Check \tg{maximum\_biomass\_density}.\\
the answer depends on $z$ & In a quasi-2D setup it must not. Usually the geometry file was written in the wrong loop order.\\
the geometry looks scrambled & The geometry file loop order. $x$ fastest, then $y$, then $z$.\\
flow never converges & \tg{tau} too close to 0.5; or \tg{delta\_P} so large the flow is not creeping; or the domain has no percolating path.\\
FBA growth is always zero & The uptake bounds are zero, or \tg{exchange\_reaction\_indices} points at the wrong reactions, or the model's objective is not the biomass reaction.\\
surrogate returns nonsense & The run has left the network's training box. Run \file{inspectSurrogate.py} and compare with your \tg{fba\_maximum\_uptake\_flux}.\\
combined type double counts growth & The kinetics file is writing \code{bioR} as well as the FBA path. It must not.\\
porosity oscillates & \tg{reopen\_fraction} too close to 1. Try 0.9.\\
\bottomrule
\end{longtable}

\section{Performance}

\begin{longtable}{@{}p{5.0cm}p{9.8cm}@{}}
\toprule
\textbf{Symptom} & \textbf{What to do}\\
\midrule\endhead
speciation dominates the run & It will. It is the most expensive part of the code by a wide margin. Turn it off unless you need it.\\
COBRApy is very slow & Expected: one to two orders of magnitude slower than GLPK per voxel. Use GLPK for production and COBRApy as a cross-check.\\
FBA dominates & Fit a surrogate. Chapter~\ref{ch:trainsurrogate}.\\
more MPI ranks made it slower & Too few voxels per rank. See Chapter~\ref{ch:cluster}.\\
the flow re-solves constantly & \tg{update\_interval} in \tg{precipitation} or \tg{dissolution} is too small, or the porosity is chattering.\\
\bottomrule
\end{longtable}

% =====================================================================
\chapter{Known limits}
\label{ch:limits}

Stated plainly, because knowing them saves more time than working around them.

\begin{itemize}[leftmargin=1.2em]
\item \textbf{Speciation is ideal.} No activity corrections, no temperature
dependence, no gas phase, no redox couples, no mineral saturation indices. Your
log\,K values are conditional constants at your own ionic strength and
25\,\textdegree C.

\item \textbf{Dissolution is not driven by saturation state.} Because there are
no solubility products, precipitation and dissolution are separate rate laws
that happen to move the same mineral, rather than one reversible reaction.

\item \textbf{Rate laws are compiled in.} Two problems with different chemistry
need two binaries.

\item \textbf{One fluid phase.} No unsaturated flow.

\item \textbf{The \file{.vti} output has unit spacing}, so ParaView shows
lattice coordinates rather than micrometres.

\item \textbf{The shipped surrogate is one organism on two substrates} over one
range. It is a demonstration of the mechanism, not a general-purpose solver
option the way GLPK is.

\item \textbf{Palabos v2.3.0 specifically.} Later versions change
data-processor interfaces.

\item \textbf{\file{src/complab3d\_glpkcpp.hh} derives from GLPKMEX} / the
Octave-Forge \file{glpkcc.cpp}, which is GPL-licensed. Its provenance is
recorded in the file header; anyone redistributing should satisfy themselves
this is correct for their situation.
\end{itemize}

% =====================================================================
\chapter{Where everything lives}
\label{ch:filemap}

\begin{longtable}{@{}p{6.0cm}p{8.8cm}@{}}
\toprule
\textbf{Path} & \textbf{What it is}\\
\midrule\endhead
\multicolumn{2}{@{}l}{\textit{You will edit these}}\\
\file{CompLaB.xml} & the run\\
\file{defineKinetics.hh} & your biotic rate laws, compiled in\\
\file{defineAbioticKinetics.hh} & your abiotic rate laws and the dissolution hook\\
\file{surrogateModel.hh} & the trained surrogate network\\
\file{CMakeLists.txt} & the Palabos path, lines 85 to 93\\
\midrule
\multicolumn{2}{@{}l}{\textit{Reference and tools}}\\
\file{CompLaB\_reference\_template.xml} & every tag, documented inline\\
\file{CompLaB\_example\_uranium\_zhao.xml} & a full 95-species production case\\
\file{extractMM.py} & genome-scale model $\rightarrow$ matrix file\\
\file{comp.sh} & Slurm template\\
\midrule
\multicolumn{2}{@{}l}{\textit{The test suite}}\\
\file{examples/} & fifteen cases, one per capability\\
\file{examples/makeExamples.py} & regenerates them\\
\file{examples/validateExamples.py} & structural check, one second\\
\file{scripts/setup\_case.sh} & assembles one case, ready to build\\
\file{tests/check\_repo.sh} & structural check over the whole tree\\
\file{examples/runExamples.slurm} & the same, under Slurm\\
\midrule
\multicolumn{2}{@{}l}{\textit{Surrogate training}}\\
\file{surrogate\_training/generateTrainingData.py} & offline FBA sweep\\
\file{surrogate\_training/trainSurrogate.py} & fit and export, no licence needed\\
\file{surrogate\_training/trainSurrogate.m} & the MATLAB path\\
\file{surrogate\_training/verifyExport.py} & compiles the export and checks it\\
\file{surrogate\_training/inspectSurrogate.py} & recovers a network's valid range\\
\midrule
\multicolumn{2}{@{}l}{\textit{Source you are unlikely to touch}}\\
\file{src/complab.cpp} & the driver: parse, set up lattices, time loop\\
\file{src/complab\_functions.hh} & XML parsing and lattice setup\\
\file{src/complab3d\_processors*.hh} & the data processors, one file per concern\\
\file{src/complab3d\_metabolic.hh} & the optional metabolic layer\\
\file{src/complab3d\_glpkcpp.hh} & the GLPK wrapper\\
\file{src/complab3d\_pythonAPI.hh} & the COBRApy bridge, C\texttt{++} side\\
\file{src/complab3d\_cobrapy.py} & the COBRApy bridge, Python side\\
\file{src/precipitationVOP.hh} & pore clogging\\
\file{src/dissolutionVOP.hh} & pore reopening\\
\bottomrule
\end{longtable}

\vspace{14pt}
\begin{cbnote}
\textbf{\color{cbTeal}Reporting a problem.}\quad
Open an issue at
\url{https://github.com/shahram444/CompLB3D-lattice-Boltzmann-FBA-surrogate-COBRApy-GLPK}
with your \file{CompLaB.xml}, the domain size, the cmake line and the terminal
output. The start-up banner is verbose and usually contains the answer. If it
is a crash, the smallest input that still reproduces it is worth more than a
large one.
\end{cbnote}
